\chapter{Assessment of the existing tools}
\label{ch:assessment}

Two existing tools in \texttt{PeleAnalysis} bear on this work.
\texttt{isosurface.cpp} extracts a multilevel isosurface by marching cubes;
\texttt{amrToFE.cpp} converts an AMR hierarchy into a finite-element mesh. The
first is the natural starting point and turns out to be the wrong one; the
second contains the construction that is actually needed.

\section{\texttt{isosurface.cpp}: a good front end, the wrong engine}

The plotfile front end is directly reusable and is the tedious part to get
right: \texttt{PlotFileData} access, component selection by name, per-level
\texttt{Geometry} carrying an explicit periodicity, \texttt{FillPatchTwoLevels}
for level-to-level data, and the fine-covered mask with its periodic-shift
pass. Roughly the first 250 lines of \texttt{main()} survive.

The engine does not, for reasons that are structural rather than qualitative.

\subsection{One isovalue is not a field topology}

Marching cubes yields a triangle soup at a single isovalue with no domain
connectivity attached. Every contour-tree and Reeb-graph
algorithm~\cite{carr2003contourtree,tierny2009loopsurgery,gueunet2016contourforests}
consumes a cell complex with vertex scalars and explicit adjacency, sweeping
sorted vertices with union--find. \texttt{Polygonise} discards precisely that:
its output is triangles of iterators into a map keyed on edges. There is no
path from that representation to critical points and their pairing.

\subsection{The coarse--fine seam is only approximately closed}

Coarse dual cells touching a covered cell emit nothing, and the fine level
covers them through its ghost band. Coarse and fine nodes on the interface
plane then land at different positions and are welded only if they compare
close. For an area integral the resulting cracks are a percent-level error; for
topology they are fatal, since a crack splits one component into two and
invents both a feature and a critical point.

\subsection{Node welding is not a strict weak ordering}

The node comparator short-circuits on an $\epsilon$ test before falling through
to a lexicographic comparison. With $a \sim b$ and $b \sim c$ but
$a \not\sim c$, transitivity fails and the behaviour of the ordered container
becomes insertion-order dependent. At $\epsilon = 10^{-15}$ this is effectively
exact comparison and largely harmless; any attempt to widen the tolerance in
order to weld coarse nodes to fine ones converts it into non-deterministic
topology.

\subsection{Piecewise-constant interpolation in the ghost band}

The coarse-to-fine fill is piecewise constant, so the field a fine box sees in
its ghost region is a staircase of coarse values. The two sides of the
interface therefore see genuinely different functions, and the staircase
manufactures exactly-flat plateaus, which are degenerate for any sweep and
create spurious critical points at every interface.

\subsection{Collation to a single rank}

Nodes and elements are gathered to the I/O processor, which rebuilds the global
sets serially. This already binds for isosurfaces at scale and is disqualifying
under~\ref{req:parallel}.

\subsection{The periodic defect is structural}

The tool's own diagnostics record it: across a periodic face the crossing is
interpolated from raw coordinates rather than the minimum image, so the node
lands mid-domain---observed inflating a reported arc length by 45\%\ in one
case. The cause is that connectivity is inferred \emph{from coordinates}. In a
complex whose adjacency is topological, the defect cannot be expressed.

\section{\texttt{amrToFE.cpp}: the construction that was needed}

This tool builds a conforming multilevel mesh, and the reason it works is a
single choice: a node's identity is the pair $(\lev,\iv)$ of level and index,
never a coordinate. Two elements share a vertex because they \emph{name the
same cell}.

The coarse--fine transition follows from that. For $\lev > 0$ the ring of cells
within \texttt{ref} of the fine grids but outside them is stamped with the
\emph{coarse parent's node}: for each coarse cell $\iv$, every one of the
$r^{\dsp}$ fine positions $r\iv + \iv_t$ receives the node
$(\lev-1, \iv)$. Coarse cells under a finer grid are marked covered, and an
element is emitted only when all of its corners are valid. Interface elements
are therefore emitted at the fine level with their coarse-side corners being
the coarse level's own nodes, by identity. Because $r^{\dsp}$ fine positions
collapse onto one coarse node, those elements carry repeated indices: they are
degenerate hexahedra---pyramids, prisms and tetrahedra written in brick
form---and that degeneracy \emph{is} the stitching. It is an alternative to the
explicit stitch-cell case tables of Weber
et~al.~\cite{weber2001crackfree,weber2012parallel} that needs no case table and
imposes no restriction on refinement ratio.

What it lacked for the present purpose:

\begin{itemize}[leftmargin=*]
\item \textbf{Periodicity by duplication, not identification.} Its periodic
  option grows the domain past the face and appends shifted copies of the
  boxes. The stamped indices there lie \emph{outside} the domain and are
  therefore different nodes from their partners on the far side. The result is
  an extended slab, correct for a seamless picture and wrong for topology: the
  complex is not a torus, so it carries no $1$-cycles.
\item \textbf{No parallelism.} The node map is built per rank with a per-rank
  counter, so identifiers collide between ranks, and the gather is commented
  out. The tool is single-rank-correct only, and the serialness is in the data
  structure rather than the interface.
\item \textbf{Legacy interfaces} predating the modern AMReX API.
\end{itemize}

\Cref{ch:complex} describes the rebuild that keeps the identity scheme and
addresses each of these.
